%! Unit Launch: Study Design — Unit 1, Lesson 1.0 — Group Activity (Answer Key)
% ONE activity for the whole class — no tiers. Every group works items 1-7 in 18 minutes;
% the extension box is for groups that finish early, not a separate assignment.
\documentclass[10pt]{article}
\usepackage{saar-article}
\usepackage{saar-key}   % pulls in -boxes; do NOT also load -boxes
\newcommand{\TallMath}[1]{$\displaystyle #1\rule[-1.4em]{0pt}{3.2em}$}

\begin{document}

\pageheader{Unit 1, Lesson 1.0}{Group Activity --- Answer Key}
% NAMESTRIP: no \namedateperiod here — mirrors the blank. Teacher-only prose goes
% in the lesson plan as \begin{teachernote}[Group Activity], never in a key.

\vspace{0.15in}

\begin{headlinebox}{frost}
\small
\textbf{Data Detectives --- The Riverbend Cross-Country Meet.} The athletic department kept a
card on every runner at Saturday's meet, and a recap blog has already made a claim about it.
Work with your group. \textbf{Every answer needs a reason}, and \textbf{every conclusion has to
name the group it applies to.} ``It's words'' and ``it's numbers'' are not reasons.
\end{headlinebox}

\vspace{0.06in}

\begin{center}
\small
\renewcommand{\arraystretch}{1.15}
\begin{tabular}{@{} l c c l c @{}}
\multicolumn{5}{c}{\textbf{The Meet Card}} \\
\toprule
\textbf{Runner} & \textbf{Bib \#} & \textbf{Finish time (min)} &
\textbf{Team color} & \textbf{Medal?} \\
\midrule
Amara & 31 & 22 & Blue & Yes \\
Ben   & 18 & 26 & Gold & No  \\
Cruz  & 45 & 24 & Blue & No  \\
Dana  &  9 & 21 & Gold & Yes \\
Eli   & 12 & 27 & Blue & No  \\
\bottomrule
\end{tabular}
\end{center}

\vspace{0.06in}

\boxguard[30]
\begin{tcolorbox}[colback=white, colframe=black!40, breakable,
    title=\textbf{Part 1 --- Read the data},
    fonttitle=\bfseries, arc=2mm, left=3mm, right=3mm, top=2mm, bottom=2mm]
\small
\begin{enumerate}[leftmargin=1.6em, itemsep=6pt, topsep=4pt]

  \item How many \textbf{individuals} does the meet card describe? \ans{5}
        \quad Not counting ``Runner,'' how many \textbf{variables}? \ans{4}

  \item Classify each recorded variable, and give the reason you used.

        \begin{center}
        \renewcommand{\arraystretch}{1.4}
        \begin{tabularx}{0.95\linewidth}{@{} p{3cm} p{3cm} X @{}}
        \toprule
        \textbf{Variable} & \textbf{Type} & \textbf{Our reason} \\
        \midrule
        Bib \#           & \ans{categorical} & \ans{it identifies a runner; its mean is meaningless} \\
        Finish time      & \ans{quantitative} & \ans{it measures an amount, in minutes} \\
        Team color       & \ans{categorical} & \ans{Blue and Gold are group names} \\
        Medal?           & \ans{categorical} & \ans{Yes/No are labels, not amounts} \\
        \bottomrule
        \end{tabularx}
        \end{center}

  \item Find the mean finish time for these five runners.

        \begin{work}
          \bar{x} &= \dfrac{22+26+24+21+27}{5} \\
                  &= \dfrac{120}{5} \\
                  &= 24 \text{ minutes}
        \end{work}

  \item ``Medal?'' is categorical, so a mean is not available. Summarize it the
        other way --- what percent of these runners medaled?

        \begin{work}
          \dfrac{2}{5} &= 0.40 \\
                       &= 40\%
        \end{work}

  \item Write one sentence saying what your answer to item 3 describes. Name the
        group and the units.
        \ansline{The five Riverbend runners on the meet card had a mean finish time of 24 minutes.}

\end{enumerate}
\end{tcolorbox}

\vspace{0.08in}

\boxguard[30]
\begin{tcolorbox}[colback=white, colframe=black!40, breakable,
    title=\textbf{Part 2 --- Judge the claim},
    fonttitle=\bfseries, arc=2mm, left=3mm, right=3mm, top=2mm, bottom=2mm]
\small
\begin{enumerate}[leftmargin=1.6em, itemsep=6pt, topsep=4pt, start=6]

  \item The recap blog wrote:

        \begin{center}
        \textit{``Riverbend's mean bib number was 23 --- a low number for a varsity
        meet. This is a young, inexperienced squad.''}
        \end{center}

        \textbf{(a)} Check the blog's arithmetic.

        \begin{work}
          \bar{x} &= \dfrac{31+18+45+9+12}{5} \\
                  &= \dfrac{115}{5} \\
                  &= 23
        \end{work}

        \textbf{(b)} The arithmetic is correct. \textbf{Is the conclusion?} Use
        the type you assigned to ``Bib \#'' in item 2 to explain what is wrong.
        \ansline{No. Bib \# is categorical --- it only identifies a runner --- so its mean is not}
        \ansline{an amount of anything. No runner wears bib 23, and it says nothing about age.}

  \item The coach says: ``Bib numbers are numbers, so of course you can average
        them. Team colors are words, so of course you cannot.'' The coach reaches
        the right answer for ``Team color'' but the \emph{wrong reason} for
        ``Bib \#.'' State the rule the coach should be using instead.
        \ansline{Ask whether arithmetic on the variable means anything --- not whether it looks numeric.}

\end{enumerate}
\end{tcolorbox}

\vspace{0.08in}

\boxguard
\begin{extensionbox}
\small
\begin{enumerate}[leftmargin=1.6em, itemsep=5pt, topsep=2pt]

  \item The school newspaper prints: \textit{``Riverbend runners average 24 minutes ---
        the fastest squad in the district.''} The arithmetic behind the $24$ is correct.
        Name \emph{two} separate things this sentence claims that the meet card cannot
        support.
        \ansline{(1) The 24 minutes describes only these five runners, not every Riverbend runner.}
        \ansline{(2) No other district team was measured, so ``fastest in the district'' has no data.}

  \item Two reporters cover the same five runners. Reporter 1 records \emph{how many
        minutes} each runner took. Reporter 2 records only whether each runner finished
        \emph{under 25 minutes} --- yes or no.

        Type for Reporter 1: \ans{quantitative} \hfill
        Type for Reporter 2: \ans{categorical}

        Same runners, same race. Explain how the \emph{question asked} --- not the
        runners --- decided the type of data.
        \ansline{Reporter 1 asked for an amount; Reporter 2 asked for a yes/no label you can only count.}

\end{enumerate}
\end{extensionbox}

\end{document}
